% Part: many-valued-logic
% Chapter: three-valued-logics
% Section: lukasiewicz

\documentclass[../../../include/open-logic-section]{subfiles}

\begin{document}

\olfileid{mvl}{inf}{luk}

\olsection{لوکاشیویچ دی منطق}

\begin{editorial}
  ایہ لامتناہی قدراں والی لوکاشیویچ دی منطق بارے حصے دا نکا جیہا مسودہ اے۔
\end{editorial}

\begin{defn}\ollabel{def:lukasiewicz}
لامتناہی قدراں والی لوکاشیویچ دی منطق~$\LogLuk[\infty]$ ایس
میٹرکس نال متعین اے۔
\begin{enumerate}
  \item معیاری قضیاتی زبان $\Lang L_0$ جیہدے وچ
  $\lnot$، $\land$، $\lor$، $\lif$ نیں۔
  \item صداقتی قدراں دا سیٹ $V_\infty$۔
  \item $1$ اکّو مقرر قدر اے، یعنی $V^+ = \{1\}$۔
  \item صداقتی فنکشن ہِٹھ دِتّے نیں۔
  \begin{align*}
    \tf{\lnot}[\LogLuk](x) & = 1 - x\\
    \tf{\land}[\LogLuk](x,y) & = \min(x,y)\\
    \tf{\lor}[\LogLuk](x,y) & = \max(x,y)\\
    \tf{\lif}[\LogLuk](x,y) & = \min(1,1-(x-y)) = \begin{cases}
      1 & \text{جے } x \le y\\
      1-(x-y) & \text{ہور صورت وچ۔}
    \end{cases}
    \end{align*}
\end{enumerate}
$m$ قدراں والی لوکاشیویچ دی منطق ایسے طرح متعین اے، بس $V = V_m$ اے۔
\end{defn}

\begin{prop}
  \olref[thr][luk]{def:lukasiewicz} وچ متعین منطق $\LogLuk[3]$،
  \olref{def:lukasiewicz} وچ متعین منطق $\LogLuk[3]$ دے عین برابر اے۔
\end{prop}

\begin{proof}
  ایہ گل \olref[thr][luk]{def:lukasiewicz} وچ دِتّے رابطیاں دے
  صداقتی جدولاں نوں \olref{def:lukasiewicz} دیاں مساواتاں توں
  ملن والے جدولاں نال ملا کے دِسدی اے۔
  \begin{center}
    \begin{tabular}{c|c} 
      $\tf{\lnot}$ & \\ 
      \hline  
      $1$ & $0$ \\ 
      $1/2$ & $1/2$ \\
      $0$ & $1$ 
    \end{tabular}
    \quad
    \begin{tabular}{c|ccc} 
      $\tf{\land}[\LogLuk[3]]$ & $1$ & $1/2$ & $0$ \\ 
      \hline 
      $1$ & $1$ & $1/2$ & $0$ \\ 
      $1/2$ & $1/2$ & $1/2$ & $0$\\ 
      $0$ & $0$ & $0$ & $0$ 
    \end{tabular}
    \\[2ex]
    \begin{tabular}{c|ccc} 
      $\tf{\lor}[\LogLuk[3]]$ & $1$ & $1/2$ & $0$ \\ 
      \hline 
      $1$ & $1$ & $1$ & $1$ \\ 
      $1/2$ & $1$ & $1/2$ & $1/2$ \\
      $0$ & $1$ & $1/2$ & $0$ 
    \end{tabular}
    \quad
    \begin{tabular}{c|ccc} 
      $\tf{\lif}[\LogLuk[3]]$ & $1$ & $1/2$ & $0$ \\ 
      \hline 
      $1$ & $1$ & $1/2$ & $0$ \\ 
      $1/2$ & $1$ & $1$ & $1/2$  \\ 
      $0$ & $1$ & $1$ & $1$ 
    \end{tabular}
  \end{center}
\end{proof}

\begin{prop}\ollabel{prop:luk-infty-m}
  جے $\Gamma \Entails[\LogLuk[\infty]] !B$، تاں $\Gamma
  \Entails[\LogLuk[m]] !B$ ہر~$m \ge 2$ لئی۔
\end{prop}

\begin{proof}
  مشق۔
\end{proof}

\begin{prob}
  \olref[mvl][inf][luk]{prop:luk-infty-m} ثابت کرو۔
\end{prob}

اصل وچ اُلٹی گل وی قائم اے۔

لامتناہی قدراں والی لوکاشیویچ دی منطق سبھ توں مقبول دھندلی
منطق اے۔ دھندلی منطق بارے لکھتاں وچ شرطیہ اکثر $\lnot !A \lor !B$
وجوں متعین ہوندا اے۔ ایس توں لامتناہی قدراں والی کلینی دی مضبوط منطق
بنے گی۔

\begin{prob}
  دکھاؤ کہ $(p \lif q) \lor (q \lif p)$، ~$\LogLuk[\infty]$
  وچ ہمیشہ صادق اے۔
\end{prob}

\end{document}
